% \iffalse meta-comment
%
% Copyright (C) 2022 by F. Pantigny
% -----------------------------------
%
% This file may be distributed and/or modified under the
% conditions of the LaTeX Project Public License, either version 1.3
% of this license or (at your option) any later version.
% The latest version of this license is in:
%
% http://www.latex-project.org/lppl.txt
%
% and version 1.3 or later is part of all distributions of LaTeX
% version 2005/12/01 or later.
%
% \fi
% \iffalse
\def\myfileversion{0.6a}
\def\myfiledate{2022/09/19}
%
%
%<*batchfile> 
\begingroup
\input l3docstrip.tex
\keepsilent
\usedir{tex/latex/cascade}
\preamble

Copyright (C) 2022 by F. Pantigny
-----------------------------------

This file may be distributed and/or modified under the
conditions of the LaTeX Project Public License, either version 1.3
of this license or (at your option) any later version.
The latest version of this license is in:

http://www.latex-project.org/lppl.txt

and version 1.3 or later is part of all distributions of LaTeX
version 2005/12/01 or later.

\endpreamble
\askforoverwritefalse
\endgroup
%</batchfile>
%
%<@@=piton>
%<*driver>
\documentclass{l3doc}
\usepackage{geometry}
\geometry{left=2.8cm,right=2.8cm,top=2.5cm,bottom=2.5cm,papersize={21cm,29.7cm}}
\usepackage{fontspec}
\usepackage[dvipsnames]{xcolor}
\usepackage{caption,tabularx}
\def\emphase{\bgroup\color{RoyalPurple}\let\next=}
\fvset{commandchars=\~\#\@,formatcom=\color{gray}}
\captionsetup{labelfont = bf}
\usepackage{ragged2e}
\usepackage[escape-inside=$$]{piton} % $$

\def\interitem{\vskip 7mm plus 2 mm minus 3mm}          
\parindent 0pt
\skip\footins = 2\bigskipamount

\begin{filecontents}[noheader,overwrite]{example1.txt}
from math import pi

def arctan(x,n=10):
    """Compute the value of arctan(x)
       n is the number of terms if the sum"""
    if x < 0:
         return -arctan(-x) # recursive call
    elif x > 1: 
         return pi/2 - arctan(1/x) 
         ## (we have used that $\arctan(x)+\arctan(1/x)=\frac{\pi}{2}$ for $x>0$)\footnote{This LaTeX escape has been done by beginning the comment by \ttfamily\#\#}
    else 
         s = 0
         for k in range(n):
              s += (-1)**k/(2*k+1)*x**(2*k+1)
         return s 
\end{filecontents}

\begin{filecontents}[noheader,overwrite]{example2.txt}
def fact(n):
    if n==0:
        return 1
    else:
        $\colorbox{yellow!50}{$return n*fact(n-1)$}$
\end{filecontents}


\EnableCrossrefs

\begin{document}
\DocInput{piton.dtx}
\end{document}
%</driver>

% \fi
% \title{The package \pkg{piton}\thanks{This document corresponds to the
% version~\myfileversion\space of \pkg{piton}, at the date of~\myfiledate.}} 
% \author{F. Pantigny \\ \texttt{fpantigny@wanadoo.fr}}
%
% \maketitle
%
% \begin{abstract}
% The package \pkg{piton} provides tools to typeset Python listings with
% syntactic highlighting by using the Lua library LPEG. It requires LuaLaTeX.
% \end{abstract}
% 
% 
% \section{Presentation}
% 
%
% The package \pkg{piton} uses the Lua library LPEG\footnote{LPEG is a
% pattern-matching library for Lua, written in C, based on \emph{parsing
% expression grammars}: \url{http://www.inf.puc-rio.br/~roberto/lpeg/}} for
% parsing Python listings and typeset them with syntactic highlighting. Since it
% uses Lua code, it works with |lualatex| only (and won't work with the other
% engines: |latex|, |pdflatex| and |xelatex|). It does not use external program
% and the compilation does not require |--shell-escape|. The compilation is very
% fast since all the parsing is done by the library LPEG, written in C.
% 
% \bigskip
% Here is an exemple of code typeset by \pkg{piton}, with the environment |{Piton}|.
%
% \bigskip
% \PitonInputFile{Example1.txt}
%
%
% \section{Installation}
% 
% The package \pkg{piton} is contained in two files: |piton.sty| and
% |piton.lua| (the LaTeX file |piton.sty| loaded by |\usepackage| will load the
% Lua file |piton.lua|). Both files must be in a repertory where LaTeX will be
% able to find them. The best way is to install them in a |texmf| tree.
% 
% 
% \section{Use of the package}
%
% In order to use the package \pkg{piton}, one has only to load the package in
% its document with the standard command |\usepackage| and remember that the
% compilation must be done with |lualatex| (and no other LaTeX engine).
%
% \medskip
% The package \pkg{piton} provides three tools to typeset Python code: the
% command |\piton|, the environment |{Piton}| and the command |\PitonInputFile|.
%
% \begin{itemize}
% \item The command |\piton| should be used to typeset small pieces of code
% inside a paragraph. \emph{Caution}: That fonction takes in its argument
% \emph{verbatim}. Therefore, it cannot be used in the argument of another
% command (however, it can be used within an environment).
% \item The environment |{Piton}| should be used to typeset multi-lines code.
% 
% \item The command |\PitonInputFile| is used to insert and typeset a whole
% external file. 
% \end{itemize}
%
% \bigskip
% It's possible to compose comments in LaTeX by beginning with |##| (it's a
% ``LaTeX escape''). The characters |##| themselves won't be printed and the
% spaces after |##| are removed.
% 
% \section{Customization}
%
% \subsection{The command \textbackslash PitonOptions}
%
% The command |\PitonOptions| provides five keys: |gobble|, |auto-gobble|,
% |env-gobble|, |line-numbers| and |all-line-numbers|.
%
% \begin{itemize}
% \item The key |gobble| takes in as value a positive integer~$n$: the first $n$
% characters are discarded (before the process of hightlighning of the code) for
% each line of the environment |{Piton}|.
%
% \item Then the key |auto-gobble| is in force, the extension \pkg{piton} computes
% the minimal value $n$ of the number of consecutives space beginning each (non
% empty) line of the environment |{Piton}| and applies |gobble| with that value
% of~$n$.
%
% \item When the key |env-gobble| is in force, \pkg{piton} applies |gobble| with
% a value of $n$ equal to the number of spaces before |\end{Piton}| on the last
% line (if that line contains only spaces).
%
% \item With the key |line-numbers|, the \emph{non empty} lines are numbered in
% the environments |{Piton}| and in the listings resulting from the use of
% |\PitonInputFile|.
% 
% \item With the key |all-line-numbers|, \emph{all} the lines are numbered,
% including the empty ones.
% \end{itemize}
%
% \bigskip
%
% \begingroup
% \fvset{commandchars=\~\&\@,formatcom=\small\color{gray}}
% \begin{Verbatim}
% ~emphase&\PitonOptions{line-numbers,auto-gobble}@
% \begin{Piton}
%     from math import pi
%
%     def arctan(x,n=10):
%         """Compute the value of arctan(x)
%            n is the number of terms if the sum"""
%         if x < 0:
%              return -arctan(-x) # recursive call
%         elif x > 1: 
%              return pi/2 - arctan(1/x) 
%              ## (on a utilisé le fait que $\arctan(x)+\arctan(1/x)=\frac{\pi}{2}$ pour $x>0$)
%         else 
%              s = 0
%              for k in range(n):
%                   s += (-1)**k/(2*k+1)*x**(2*k+1)
%              return s 
% \end{Piton}
% \end{Verbatim}
% \endgroup
%
%
% \begingroup
% \PitonOptions{line-numbers,auto-gobble}
% \PitonInputFile{Example1.txt}
% \endgroup
% 
% \subsection{The option escape-inside}
%
% The option |escape-inside| must be used when loading the package \pkg{piton}
% (that is to say in the instruction |\usepackage|). For technical reasons, it
% can't be used in the command |\PitonOptions|. That option takes in as value
% two characters which will be used to delimit pieces of code which will
% composed in LaTeX.
%
% \medskip
% In the following example, we assume that the extension \pkg{piton} has been
% loaded by the following instruction.
%
% \begin{Verbatim}
% \usepackage[~emphase#escape-inside=$$@]{piton}
% \end{Verbatim}
%
% \medskip
% In the following code, which is a recursive programmation of the mathematical
% factorial, we decide to highlight in yellow the instruction which contains the
% recursive call. 
% \begin{Verbatim}
% \begin{Piton}
% def fact(n):
%     if n==0:
%         return 1
%     else:
%         ~emphase#$\colorbox{yellow!50}{$@return n*fact(n-1)~emphase#$}$@
% \end{Piton}
% \end{Verbatim}
%
% \PitonInputFile{Example2.txt}
%
% \bigskip
%
% \emph{Caution} : The escape to LaTeX allowed by the characters of
% |escape-inside| is not active in the strings nor in the Python comments
% (however, it's possible to have a whole Python comment composed in LaTeX by
% beginning it with |##|). 
%
%
% \subsection{The styles}
%
% The package \pkg{piton} provides the command |\SetPitonStyle| to customize the
% different styles used to format the syntactic elements of the Python listings.
% The customizations done by that command are limited to the current TeX group.
%
% \bigskip
% The command |\SetPitonStyle| takes in as argument a comma-separated list of
% \textsl{key=value} pairs. The keys are names of styles and the value are LaTeX
% formatting instructions.
%
% \bigskip
% These LaTeX instructions must be formatting instructions such as
% |\color{...}|, |\bfseries|, |\slshape|, etc. (the commands of this kind are
% sometimes called \emph{semi-global} commands). It's also possible to put,
% \emph{at the end of the list of instructions}, a LaTeX command taking exactly
% one argument.
% 
% \bigskip
% Here an example which changes the style used to highlight, in the definition
% of a Python function, the name of the function which is defined.
%
% \begin{verbatim}
% \SetPitonStyle 
%   { Name.Function = \bfseries \setlength{\fboxsep}{1pt}\colorbox{yellow!50} }
% \end{verbatim}
%
% In that example, |\colorbox{yellow!50}| must be considered as the name of a
% LaTeX command which takes in exactly one argument, since, usually, it is used
% with the syntax |\colorbox{yellow!50}{text}|.
%
% \medskip
% \begingroup
% \SetPitonStyle 
%   { Name.Function = \bfseries \setlength{\fboxsep}{1pt}\colorbox{yellow!50} }
% With that setting, we will have : \piton{def cube(x) : return x * x * x }
% \endgroup
% 
% \medskip
% The different styles are described in the table \ref{Semantic}.
%
%
%
% \begin{table}[b]
% \centering
% \caption{Usage of the different styles}
% \label{Semantic}
% \begin{tabular}{@{}>{\ttfamily}l>{\RaggedRight\arraybackslash}p{9cm}@{}}
% \toprule
% \normalfont Style  & Usage \\
% \midrule
% Number & the numbers \\
% String.Short & the short strings (between \texttt{'} or \verb|"|) \\
% String.Long & the long strings (between \texttt{'''} or \verb|"""|) except the
% documentation strings \\ 
% String & that keys sets both |String.Short| and |String.Long| \\
% String.Doc & the documentation strings \\
% String.Interpol & the syntactic elements of the fields of the f-strings (that
% is to say the characters \texttt{\{}, \texttt{\}} and \texttt{:}) \\
% Operator & the following operators : \texttt{!= == << >> - \~{} + / * \% = < > \& .} \verb+|+ \verb|@|\\
% Operator.Word & the following operators : \texttt{in}, \texttt{is},
% \texttt{and}, \texttt{or} et \texttt{not} \\ 
% Name.Builtin & the predefined functions of Python \\
% Name.Function & the name of the functions defined by the user, at the point of
% their definition (that is to say after the keyword |def|) \\
% Name.Decorator & the decorators (instructions beginning by \verb|@| in the classes) \\
% Name.Namespace & the name of the modules (= external libraries) \\
% Name.Class & the name of the classes at the point of their definition \\
% Exception & the names of the exceptions (eg: \texttt{SyntaxError}) \\
% Comment & the comments beginning with \verb|#| \\
% LaTeX & the comments beginning by \verb|##| which are composed in LaTeX by
% \pkg{piton} (\verb|##| is an espace sequence to LaTeX) \\
% Keyword.Constant & |True|, |False| and |None| \\
% Keyword & the following keywords :
%           \ttfamily assert, break, case, continue, del,
%           elif, else, except, exec, finally, for, from, 
%           global, if, import, lambda, non local,
%           pass, raise, return, try, while,
%           with, yield, yield from.\\
% \bottomrule 
% \end{tabular}
% \end{table}
%
%
%
% 
% \subsection{Creation of new environments}
%
% Since the environment |{Piton}| has to catch its body in a special way (more
% or less as verbatim text), it's not possible to construct new environments
% directly over the environment |{Piton}|
%
% That's why \pkg{piton} provides a command |\NewPitonEnvironment|. That
% command takes in three mandatory arguments. 

% That command has the same syntax as the classical environment
% |\NewDocumentEnvironment|.
% 
% \bigskip
% With the following instruction, a new environment |{Python}| will be
% constructed with the same behaviour as |{Piton}|:
%
% \verb|\NewPitonEnvironment{Python}{}{}{}|
% 
% \bigskip
% If on wished to format Python code in a box in a box of \pkg{tcolorbox}, it's
% possible to define an environment |{Python}| with the following code:
%
%\begin{verbatim}
% \NewPitonEnvironment{Python}{}
%   {\begin{tcolorbox}}
%   {\end{tcolorbox}}
% \end{verbatim}
%
%
% \clearpage
%
% \section{Implementation}
%
%     \begin{macrocode}
\NeedsTeXFormat{LaTeX2e}
\RequirePackage{l3keys2e}
\ProvidesExplPackage
  {piton}
  {\myfiledate}
  {\myfileversion}
  {Highlight Python codes with LPEG on LuaLaTeX}
%    \end{macrocode}
%
% \bigskip
%    \begin{macrocode}
\msg_new:nnn { piton } { LuaLaTeX~mandatory }
  { The~package~'piton'~must~be~used~with~LuaLaTeX.\\ It~won't~be~loaded. }
\sys_if_engine_luatex:F { \msg_critical:nn { piton } { LuaLaTeX~mandatory } }
%    \end{macrocode}
% 
% \bigskip
%    \begin{macrocode}
\RequirePackage { luatexbase } 
%    \end{macrocode}
%
% \bigskip
% We define a set of keys |piton/package| for these options.
%    \begin{macrocode}
\keys_define:nn { piton / package }
  { 
    escape-inside .tl_set:N = \c_@@_escape_inside_tl ,
    unknown .code:n = \msg_error:nn { piton } { unknown~key~for~package }
  }
%    \end{macrocode}
%
%    \begin{macrocode}
\msg_new:nnn { piton } { unknown~key~for~package }
  {
    Unknown~key\\
    You~have~used~the~key~'\l_keys_key_str'~but~the~only~key~available~here~
    is~the~key~'escape-inside'.\\
    That~key~will~be~ignored.
  }
%    \end{macrocode}
% 
% \bigskip
%    \begin{macrocode}
\tl_clear_new:N \c_@@_escape_inside_tl
%    \end{macrocode}
%
% We process the options when the package is loaded (with |\usepackage|).
%    \begin{macrocode}
\ProcessKeysOptions { piton / package }
%    \end{macrocode}
%
% \bigskip
%    \begin{macrocode}
\begingroup
\cs_new_protected:Npn \@@_set_escape_char:nn #1 #2 
  {
    \directlua { begin_escape = "#1" }
    \directlua { end_escape = "#2" }
  }
\cs_generate_variant:Nn \@@_set_escape_char:nn { x x }
\@@_set_escape_char:xx 
  { \tl_head:V \c_@@_escape_inside_tl }
  { \tl_tail:V \c_@@_escape_inside_tl }
\endgroup
%    \end{macrocode}
% 
% 
% \bigskip
%    \begin{macrocode}
\AtBeginDocument
  {
    \@ifpackageloaded { xcolor }
      { }
      { \msg_fatal:nn { piton } { xcolor~not~loaded } }
  }
%    \end{macrocode}
%
%    \begin{macrocode}     
\msg_new:nnn { piton } { xcolor~not~loaded }
  {
    xcolor~not~loaded \\
    The~package~'xcolor'~is~required~by~'piton'.\\
    This~error~is~fatal.
  }
%    \end{macrocode}
% 
% \bigskip
%    \begin{macrocode}
\msg_new:nnn { piton } { piton.lua~not~found }
  { 
    The~file~'piton.lua'~can't~be~found.\\ 
    The package~'piton'~won't be loaded.
  }
%    \end{macrocode}
%
% \bigskip
%    \begin{macrocode}
\file_if_exist:nF { piton.lua }
  { \msg_critical:nn { piton } { piton.lua~not~found} }
%    \end{macrocode}
% 
%    \begin{macrocode}
\lua_now:e { require("piton.lua") }
%    \end{macrocode}
% 
% \bigskip
% The following function has not a name with the conventions of L3 because it
% will be used in the Lua code.
%    \begin{macrocode}
\cs_new:Npn \pitonEOL 
  { 
    \par \leavevmode 
    \@@_print_number:
  }
%    \end{macrocode}
%
% \medskip
%    \begin{macrocode}
\cs_new_protected:Npn \@@_print_number:
  { 
    \bool_if:NT \l_@@_line_numbers_bool 
      { \@@_actually_print_number: }
  }
%    \end{macrocode}
%
% \medskip
%    \begin{macrocode}
\cs_new_protected:Npn \@@_actually_print_number:
  {
%    \end{macrocode}
% The empty lines in the code are not numbered.
%    \begin{macrocode}
    \bool_if:NF \l_@@_all_line_numbers_bool
      { \peek_meaning:NF \pitonEOL }
    \@@_actually_print_number_i: 
  }
%    \end{macrocode}
%
%    \begin{macrocode}
\cs_new_protected:Npn \@@_actually_print_number_i:
  {
    \int_incr:N \l_@@_lineno_int 
    \hbox_overlap_left:n
      { 
        { \color{gray} \footnotesize \int_to_arabic:n \l_@@_lineno_int }
        \quad 
      }
  }
%    \end{macrocode}
% 
% \bigskip
% The following counter will be used to count the lines in the code when the
% user requires the numbers of the lines to be printed.
%    \begin{macrocode}
\int_new:N \l_@@_lineno_int 
%    \end{macrocode}
% 
% \bigskip
% The following integer is the number of characters to gobble on the left side
% of the Python listings. Of course, the initial value is $0$.
%    \begin{macrocode}
\int_new:N \l_@@_gobble_int 
%    \end{macrocode}
%
%
% \bigskip
%    \begin{macrocode}
\cs_new_protected:Npn \@@_define_gobble_syntax:n #1
  { \lua_now:n { define_gobble_syntax(#1) } }
%    \end{macrocode}
% 
% \bigskip
%    \begin{macrocode}
\NewDocumentCommand { \piton } { v }
  {
    \group_begin:
    \ttfamily
    \lua_now:e { Parse(token.scan_argument()) } { #1 }
    \group_end:
  }
%    \end{macrocode}
% 
% \bigskip
%    \begin{macrocode}
\NewDocumentCommand { \PitonInputFile } { m }
  {
    \group_begin:
      \ttfamily
      \bool_if:NT \l_@@_line_numbers_bool 
        { 
          \@@_actually_print_number:
          \vspace{-\baselineskip}
        }
      \lua_now:e { ParseFile(token.scan_argument()) } { #1 }
    \group_end:
  }
%    \end{macrocode}
%
% \subsection{PitonOptions}
%
%    \begin{macrocode}
\bool_new:N \l_@@_line_numbers_bool 
\bool_new:N \l_@@_all_line_numbers_bool 
%    \end{macrocode}
% 
%    \begin{macrocode}
\keys_define:nn { PitonOptions }
  {
    gobble           .int_set:N         = \l_@@_gobble_int , 
    gobble           .value_required:n  = true ,
    auto-gobble      .code:n            = \int_set:Nn \l_@@_gobble_int { -1 } , 
    auto-gobble      .value_forbidden:n = true ,
    env-gobble       .code:n            = \int_set:Nn \l_@@_gobble_int { -2 } , 
    env-gobble       .value_forbidden:n = true ,
    line-numbers     .bool_set:N        = \l_@@_line_numbers_bool ,
    line-numbers     .default:n         = true ,
    all-line-numbers .code:n = 
      \bool_set_true:N \l_@@_line_numbers_bool 
      \bool_set_true:N \l_@@_all_line_numbers_bool ,
    all-line-numbers .value_forbidden:n = true  ,
    unknown          .code:n = 
      \msg_error:nn { piton } { Unknown~key~for~PitonOptions }
  }
%    \end{macrocode}
%
%  \bigskip
%    \begin{macrocode}
\msg_new:nnn { piton } { Unknown~key~for~PitonOptions }
  { 
    Unknown~key \\
    The~key~'\l_keys_key_str'~is~unknown~for~\token_to_str:N \PitonOptions.~The~
    available~keys~are:~all-line-numbers,~auto-gobble,~env-gobble,~gobble~and~
    line-numbers.\\
    If~you~go~on,~that~key~will~be~ignored.
  }
%    \end{macrocode}
% 
% \bigskip
% The argument of |\PitonOptions| is provided by currification.
%    \begin{macrocode}
\NewDocumentCommand \PitonOptions { }
  { \keys_set:nn { PitonOptions } }
%    \end{macrocode}
% 
% \bigskip 
%    \begin{macrocode}
\NewDocumentCommand { \NewPitonEnvironment } { m m m m }
  {
%    \end{macrocode}
% We construct a TeX macro which will catch as its argument all the tokens
% until \textsl{\texttt{newline}} + |\end{Piton}| with, in that
% \textsl{\texttt{newline}}| + \end{Piton}|, the catcode of
% \textsl{\texttt{newline}}, |\|, |{| and |}| equal to 12 (``\texttt{other}'').
% The latter explains why the definition of that function is a bit complicated.
%    \begin{macrocode}
    \use:x 
      { 
        \cs_set_protected:Npn 
          \use:c { __piton_collect_ #1 :w } 
          ####1 
          \c_backslash_str end \c_left_brace_str #1 \c_right_brace_str
      }
          {
            \group_end:
            \par \addvspace { 0.5 em }
            {
              \dim_set_eq:NN \parindent \c_zero_dim 
              \ttfamily
              \bool_if:NT \l_@@_line_numbers_bool 
                { 
                  \@@_actually_print_number:
                  \vspace{-\baselineskip}
                }
              \int_case:nnF \l_@@_gobble_int 
                { 
                  0 
%    \end{macrocode}
% Be careful: the last argument is provided by currification.
%    \begin{macrocode}
                  { \lua_now:e { Parse(token.scan_argument()) } }
                  { -1 }
                  { \lua_now:e { AutoGobbleParse(token.scan_argument()) } } 
                  { -2 }
                  { \lua_now:e { EnvGobbleParse(token.scan_argument()) } } 
                }
                {
                   \exp_args:NV \@@_define_gobble_syntax:n \l_@@_gobble_int 
                   \lua_now:e { GobbleParse(token.scan_argument()) } 
                }
              { ##1 }
            }
            \par \addvspace { 0.5 em } 
%    \end{macrocode}
% The following |\end{#1}| is only for the groups and the stack of
% environments of LaTeX.
%    \begin{macrocode}
            \end { #1 }
          }   
%    \end{macrocode}
% 
% 
% \bigskip
% We can now define the new environment.
%    \begin{macrocode}
    \NewDocumentEnvironment { #1 } { #2 } 
      {
        #3 
        \group_begin:
        \tl_map_function:nN 
          { \ \\ \{ \} \$ \& \# \^ \_ \% \~ } 
          \char_set_catcode_other:N 
        \use:c { __piton_collect_ #1 :w }
      }
      { #4 }
%    \end{macrocode}
%
%    \begin{macrocode}
    \AddToHook { env / #1 / begin } { \char_set_catcode_other:N \^^M }
  }
%    \end{macrocode}
%
% \bigskip
%    \begin{macrocode}
\NewPitonEnvironment { Piton } { } { } { }
%    \end{macrocode}
% 
% \bigskip
%
%    \begin{macrocode}
\NewDocumentCommand { \PitonStyle } { m } { \csname pitonStyle#1\endcsname }
%    \end{macrocode}
%
% \bigskip
%    \begin{macrocode}
\NewDocumentCommand { \SetPitonStyle } { } { \keys_set:nn { piton } }
%    \end{macrocode}
% 
%    \begin{macrocode}
\cs_new_protected:Npn \@@_math_scantokens:n #1 
  { \normalfont \scantextokens { $#1$ } }
%    \end{macrocode}
% 
% \bigskip
%    \begin{macrocode}
\keys_define:nn { piton }
  {
    String.Interpol  .tl_set:c = pitonStyle String.Interpol , 
    String.Interpol  .value_required:n = true ,
    FormattingType   .tl_set:c = pitonStyle FormattingType ,
    FormattingType   .value_required:n = true ,
    Dict.Value       .tl_set:c = pitonStyle Dict.Value ,
    Dict.Value       .value_required:n = true ,
    Name.Decorator   .tl_set:c = pitonStyle Name.Decorator ,
    Name.Decorator   .value_required:n = true ,
    Name.Function    .tl_set:c = pitonStyle Name.Function ,
    Name.Function    .value_required:n = true ,
    Keyword          .tl_set:c = pitonStyle Keyword ,
    Keyword          .value_required:n = true ,
    Keyword.Constant .tl_set:c = pitonStyle Keyword.Constant ,
    Keyword.constant .value_required:n = true ,
    String.Doc       .tl_set:c = pitonStyle String.Doc ,
    String.Doc       .value_required:n = true ,
    Interpol.Inside  .tl_set:c = pitonStyle Interpol.Inside ,
    Interpol.Inside  .value_required:n = true ,
    String.Long      .tl_set:c = pitonStyle String.Long ,
    String.Long      .value_required:n = true ,
    String.Short     .tl_set:c = pitonStyle String.Short ,
    String.Short     .value_required:n = true ,
    String           .meta:n = { String.Long = #1 , String.Short = #1 } ,
    Comment.Math     .tl_set:c = pitonStyle Comment.Math , 
    Comment.Math     .default:n = \@@_math_scantokens:n ,
    Comment.Math     .initial:n = ,
    Comment          .tl_set:c = pitonStyle Comment , 
    Comment          .value_required:n = true ,
    InitialValues    .tl_set:c = pitonStyle InitialValues ,
    InitialValues    .value_required:n = true ,
    Number           .tl_set:c = pitonStyle Number , 
    Number           .value_required:n = true ,
    Name.Namespace   .tl_set:c = pitonStyle Name.Namespace , 
    Name.Namespace   .value_required:n = true ,
    Name.Class       .tl_set:c = pitonStyle Name.Class , 
    Name.Class       .value_required:n = true ,
    Name.Builtin     .tl_set:c = pitonStyle Name.Builtin , 
    Name.Builtin     .value_required:n = true ,
    Name.Type        .tl_set:c = pitonStyle Name.Type , 
    Name.Type        .value_required:n = true ,
    Operator         .tl_set:c = pitonStyle Operator , 
    Operator         .value_required:n = true ,
    Operator.Word    .tl_set:c = pitonStyle Operator.Word , 
    Operator.Word    .value_required:n = true ,
    Post.Function    .tl_set:c = pitonStyle Post.Function , 
    Post.Function    .value_required:n = true ,
    Exception        .tl_set:c = pitonStyle Exception ,
    Exception        .value_required:n = true ,
    Comment.LaTeX    .tl_set:c = pitonStyle Comment.LaTeX ,
    Comment.LaTeX    .value_required:n = true ,
    unknown          .code:n = \msg_error:nn { piton }{ Unknown~key~for~SetPitonStyle }
  }
%    \end{macrocode}
% 
% 
% \bigskip
%    \begin{macrocode}
\msg_new:nnn { piton } { Unknown~key~for~SetPitonStyle } {
  The~style~'\l_keys_key_str'~is~unknown.\\
  This~key~will~be~ignored.\\
  The~available~styles~are~(in~alphabetic~order):~
  Comment,~
  Comment.LaTeX,~
  Dict.Value,~
  Exception,~
  InitialValues,~ 
  Keyword,~
  Keyword.Constant,~
  Name.Builtin,~
  Name.Class,~
  Name.Decorator,~
  Name.Function,~
  Name.Namespace,~
  Number,~
  Operator,~
  Operator.Word,~
  String,~
  String.Doc,~
  String.Long,~
  String.Short,~and~
  String.Interpol. }
%    \end{macrocode}
%
% \bigskip
%    \begin{macrocode}
\SetPitonStyle
  {                                                       
    Comment          = \color[HTML]{0099FF} \itshape , 
    Exception        = \color[HTML]{CC0000} ,
    Keyword          = \color[HTML]{006699} \bfseries ,               
    Keyword.Constant = \color[HTML]{006699} \bfseries ,               
    Name.Builtin     = \color[HTML]{336666} ,
    Name.Decorator   = \color[HTML]{9999FF}, 
    Name.Class       = \color[HTML]{00AA88} \bfseries ,
    Name.Function    = \color[HTML]{CC00FF} , 
    Name.Namespace   = \color[HTML]{00CCFF} , 
    Number           = \color[HTML]{FF6600} ,
    Operator         = \color[HTML]{555555} ,
    Operator.Word    = \bfseries ,
    String           = \color[HTML]{CC3300} ,         
    String.Doc       = \color[HTML]{CC3300} \itshape , 
    String.Interpol  = \color[HTML]{AA0000} ,
    Comment.LaTeX    = \normalfont \color[rgb]{.468,.532,.6} , 
    Name.Type        = \color[HTML]{336666} ,
    InitialValues    = \piton ,
    Dict.Value       = \piton ,
    Post.Function    = \piton ,
    Interpol.Inside  = \color{black}\piton ,
  }
%    \end{macrocode}
%
%
% \vspace{1cm}
% \tableofcontents
% \end{document}
%
% Local Variables:
% TeX-fold-mode: t
% TeX-fold-preserve-comments: nil
% flyspell-mode: nil
% fill-column: 80
% End:






